一张扣着的扑克牌是 A 的概率是 \(1/13\)。有人瞄了一眼告诉你``它是黑桃''，这个数仍然是 \(1/13\)；如果那句话换成``它是红色的''，这个数就变成了 \(0\)。同一张牌，两句话，两个答案，没有哪一个更正确——它们回答的是不同信息状态下的不同问题。

概率不是贴在事物身上的标签，而是相对于某个信息状态的一笔账。这一章要处理的共同困难只有一句：新信息到手之后，原来那笔账该怎么改。

四篇共用同一个动作，截断再归一——把样本空间缩到已知信息之内，再把剩下的部分整体放大，使它自己的总概率回到 \(1\)。这个动作立刻带出一条警告：\(P(A\given B)\) 与 \(P(B\given A)\) 截出的是两个不同的世界，数值上一般毫不相干。检测的技术指标属于前者，你拿到报告后真正想知道的属于后者。把两者对调，在法庭上叫检察官谬误，在监控台前则表现为对告警的系统性误读。

\subsection{本章篇目}

\begin{itemize}
\tightlist
\item \hyperref[sec:05-Probability/02-Conditional-Probability/01]{``条件概率与乘法法则''}：条件化换掉的不是公式而是全集；链式法则能把任意长的过程拆成逐步更新，并且不依赖任何独立性假设。
\item \hyperref[sec:05-Probability/02-Conditional-Probability/02]{``全概率公式与 Bayes 更新''}：后验赔率等于先验赔率乘以似然比。基础率极低时，即使检测很准，阳性结果也大概率是假阳性。
\item \hyperref[sec:05-Probability/02-Conditional-Probability/03]{``独立性：从两事件到多个事件''}：独立是概率分配的性质，不是集合几何的性质；多个事件的独立要求所有子集都按乘积分解，两两验过并不算数。
\item \hyperref[sec:05-Probability/02-Conditional-Probability/04]{``条件独立与信息结构''}：独立与条件独立互不蕴含。给定共同原因会让相关性消失，给定共同结果会让相关性凭空出现。
\end{itemize}

读完全章只需要\hyperref[sec:05-Probability/01-Probability-Spaces/02]{``概率公理与概率空间''}里的归一化与可加性两条。第四篇是概率图模型与因果推断的入口，值得比其余三篇多花些时间。

给自己留一个习惯：每当一个概率数字发生变化，先说清它是相对于什么信息而言的。这个习惯比本章任何一条公式都耐用。
